% R08 v1.6：已核验的读者延伸阅读。只写实际核验到的书目信息，不写未核验页码。
\begin{enumerate}[leftmargin=2.4em,label={[\arabic{enumi}]}]
  \item 李航：《统计学习方法（第2版）》，清华大学出版社，2019。\par
  \emph{用途：}本讲义监督学习、无监督学习、采样方法、主题模型与\PageRank{}等经典主题的范围基准；本讲义的组织、推导、例题和练习表述均为重新编排，不暗示原作者或出版社参与编写。

  \item Sheldon Axler, \emph{Linear Algebra Done Right}, 3rd ed., Springer, 2015. DOI: \href{https://doi.org/10.1007/978-3-319-11080-6}{10.1007/978-3-319-11080-6}.\par
  \emph{用途：}有限维向量空间、线性映射、特征结构、内积空间与谱定理的标准背景。

  \item Thomas M. Cover and Joy A. Thomas, \emph{Elements of Information Theory}, 2nd ed., Wiley, 2006. DOI: \href{https://doi.org/10.1002/047174882X}{10.1002/047174882X}.\par
  \emph{用途：}熵、相对熵、互信息及其等号条件的标准来源。

  \item Stephen Boyd and Lieven Vandenberghe, \emph{Convex Optimization}, Cambridge University Press, 2004. \href{https://web.stanford.edu/~boyd/cvxbook/}{作者公开页面}.\par
  \emph{用途：}凸集、凸函数、拉格朗日对偶、KKT条件与凸优化建模。

  \item Jorge Nocedal and Stephen J. Wright, \emph{Numerical Optimization}, 2nd ed., Springer, 2006. DOI: \href{https://doi.org/10.1007/978-0-387-40065-5}{10.1007/978-0-387-40065-5}.\par
  \emph{用途：}线搜索、拟牛顿法、数值停止准则以及优化算法的条件与复杂度口径。

  \item Kevin P. Murphy, \emph{Probabilistic Machine Learning: An Introduction}, MIT Press, 2022. \href{https://probml.github.io/book1}{作者公开页面}.\par
  \emph{用途：}概率建模、监督/无监督学习、潜变量模型与近似推断的现代统一背景。

  \item J. R. Norris, \emph{Markov Chains}, Cambridge University Press, 1997.\par
  DOI: \href{https://doi.org/10.1017/CBO9780511810633}{10.1017/CBO9780511810633}.\par
  \emph{用途：}离散时间马尔可夫链、平稳分布、不可约性、周期性与遍历结论的严格条件。

  \item Christian P. Robert and George Casella, \emph{Monte Carlo Statistical Methods}, 2nd ed., Springer, 2004. DOI: \href{https://doi.org/10.1007/978-1-4757-4145-2}{10.1007/978-1-4757-4145-2}.\par
  \emph{用途：}随机数生成、Monte Carlo积分、\MetropolisHastings{}、Gibbs抽样与收敛诊断。

  \item Daniel D. Lee and H. Sebastian Seung, ``Learning the Parts of Objects by Non-negative Matrix Factorization,'' \emph{Nature}, 401:788--791, 1999. DOI: \href{https://doi.org/10.1038/44565}{10.1038/44565}.\par
  \emph{用途：}非负矩阵分解的非负约束、加性表示与经典乘法更新背景。

  \item David M. Blei, Andrew Y. Ng, and Michael I. Jordan, ``Latent Dirichlet Allocation,'' \emph{Journal of Machine Learning Research}, 3:993--1022, 2003. \href{https://www.jmlr.org/papers/v3/blei03a.html}{JMLR论文页}.\par
  \emph{用途：}LDA生成模型、变分推断和经验Bayes参数估计的原始论文来源。

  \item John Lafferty, Andrew McCallum, and Fernando Pereira, ``Conditional Random Fields: Probabilistic Models for Segmenting and Labeling Sequence Data,'' \emph{Proceedings of ICML}, 2001. \href{https://repository.upenn.edu/bitstreams/4905e2c0-e9d5-4961-804b-973de8bdfc7c/download}{机构存档}.\par
  \emph{用途：}条件随机场的条件建模、特征函数和序列标注背景。

  \item Lawrence Page, Sergey Brin, Rajeev Motwani, and Terry Winograd, ``The \PageRank{}\ Citation Ranking: Bringing Order to the Web,'' Technical Report, Stanford InfoLab, 1999. \href{https://ilpubs.stanford.edu/422/}{Stanford存档}.\par
  \emph{用途：}\PageRank{}的随机浏览解释、链接矩阵与幂迭代背景。
\end{enumerate}
